\clearpage
\SourceLeaf{34}

\begin{center}\large BEWEIS DER IRRATIONALITÄT DER TANGENTEN RATIONALER BÖGEN IN EINER NEUEN GESTALT.\end{center}

\noindent\hspace*{\parindent}Die unendlichen Reihen

\[
P = 1-\frac{1}{2} . \frac{mm}{nn}+\frac{1}{2.4} . \frac{1}{1.3} . \frac{m^4}{n^4}-\frac{1}{2.4.6} . \frac{1}{1.3.5} . \frac{m^6}{n^6}+\text{und\,so\,weiter}
\]
\[
\begin{aligned}
P_1 &= \frac{m}{n}-\frac{1}{2} . \frac{1}{1.3} . \frac{m^2}{n^3}\\ &\quad+\frac{1}{2.4} . \frac{1}{1.3.5} . \frac{m^5}{n^5}-\frac{1}{2.4.6} . \frac{1}{1.3.5.7} . \frac{m^7}{n^7}+\text{und\,so\,weiter}
\end{aligned}
\]
\[
\begin{aligned}
P_2 &= \frac{1}{1.3} . \frac{m^3}{n.n}-\frac{1}{2} . \frac{1}{1.3.5} . \frac{m^5}{n^4}\\ &\quad+\frac{1}{2.4} . \frac{1}{1.3.5.7} . \frac{m^7}{n^6}-\frac{1}{2.4.6} . \frac{1}{1.3.5.7.9} . \frac{m^9}{n^8}+\text{und\,so\,weiter}
\end{aligned}
\]
\[
\begin{aligned}
P_3 &= \frac{1}{1.3.5} . \frac{m^5}{n^3}-\frac{1}{2} . \frac{1}{1.3.5.7} . \frac{m^7}{n^5}\\ &\quad+\frac{1}{2.4} . \frac{1}{1.3.5.7.9} . \frac{m^9}{n^7}-\frac{1}{2.4.6} . \frac{1}{1.3.5.7.9.11} . \frac{m^{11}}{n^9}+\text{und\,so\,weiter}
\end{aligned}
\]

\noindent und allgemein, von $\theta = 1$ an,

\[
\begin{aligned}
P_\theta = {}& \frac{1}{1.3.5\ldots(2\theta-1)} . \frac{m^{2\theta-1}}{n^\theta}
-\frac{1}{2} . \frac{1}{1.3.5\ldots(2\theta+1)} . \frac{m^{2\theta+1}}{n^{\theta+2}}
\\ &\,+\frac{1}{2.4} . \frac{1}{1.3.5\ldots(2\theta+3)} . \frac{m^{2\theta+3}}{n^{\theta+4}}
\\ &\,-\frac{1}{2.4.6} . \frac{1}{1.3.5\ldots(2\theta+5)} . \frac{m^{2\theta+5}}{n^{\theta+6}}+\text{und\,so\,weiter},
\end{aligned}
\]

\noindent wobei $m$, $n$ zunächst beliebige gegebene positive Größen bedeuten,
sind offenbar alle spätestens nach hinreichend weiter Fortsetzung
konvergent, und zwar schneller als jede fallende geometrische Folge.
Jede hat daher einen bestimmten endlichen Zahlenwert.

Ist bereits das zweite Glied von $P_\theta$ kleiner als das erste oder
gleich diesem, so ist das dritte von selbst kleiner als das zweite, das
vierte kleiner als das dritte und\,so\,weiter. Hierfür genügt es,
dass $2\theta+1$ mindestens $\dfrac{mm}{2nn}$ ist. Der
Wert von $P_\theta$ ist dann jedenfalls positiv, kleiner als das erste Glied und
größer als die Summe der beiden ersten Glieder.
